\centerline{\bf \Huge Домашнее задание.}
\centerline{\bf \Huge Признаки равенства треугольников.}

\begin{flushright}
        \scriptsize{Нет царского пути в геометрии\\Евклид.}
\end{flushright}

\problem{1}На сторонах $AC$ и $BC$ равнобедренного треугольника $ABC$ с основанием $AB$, отмечены точки $A_1$ и $B_1$ так, что $CA_1 = CB_1$.
Докажите что $\triangle CAB_1 = \triangle CBA_1$ и $\triangle ABB_1 = \triangle BAA_1$.

\problem{2}$ABCD$ – прямоугольник, точка $M$ делит $BC$ пополам. Найдите стороны прямоугольника, если его периметр 24 и $MA \perp MD$.

\problem{3}Вне прямоугольного треугольника $ABC$ на его катетах $AC$ и $BC$ построены квадраты $ACDE$ и $BCFG$. Продолжение медианы CM треугольника ABC пересекает прямую $DF$ в точке $N$. Докажите, что $CN$ высота треугольника $CDF$.

\problem{4}На стороне $AC$ равнобедренного треугольника $ABC$ отмечена точка $M$. Из неё покатился шарик в сторону середины $X$ стороны $AC$ и отскочил от неё по закону "угол падения равен углу отражения". После отскока шарик попал на сторону $BC$ в точку $Z$. Докажите, что $ZB = PC$.

\problem{5}Во сколько раз одна сторона прямоугольника больше другой?

\begin{figure}[h] 
    \includegraphics[scale=1.5]{Осень_2023/Киви/geom_1.png}
\end{figure}

\problem{$6^*$}
\noindentа) На гипотенузе $BC$ прямоугольного треугольника $ABC$ выбрана точка $K$ так, что  $AB = AK$.  Отрезок $AK$ пересекает биссектрису $CL$ в её середине. Найдите острые углы треугольника $ABC$.

\noindentб) Точки $M$ и $N$ – середины боковых сторон $AB$ и $CD$ трапеции $ABCD$. Перпендикуляр, опущенный из точки $M$ на диагональ $AC$, и перпендикуляр, опущенный из точки $N$ на диагональ $BD$, пересекаются в точке $P$. Докажите, что  $PA = PD$.